\section{挑战与展望}

\subsection{现存技术瓶颈}

尽管ISAC技术发展迅速，仍面临以下关键挑战：

\textbf{（1）自干扰消除的工程化难题。}100\,dB以上的宽带自干扰消除在工程实现中面临多重困难：功率放大器（PA）的AM/AM和AM/PM非线性效应导致数字域残余干扰难以精确建模；宽带系统中频率选择性自干扰信道的实时跟踪增加了计算复杂度；温度漂移和器件老化导致射频域对消电路需要频繁校准。

\textbf{（2）通信与感知的性能权衡。}有限的时频资源和发射功率需要在通信QoS与感知探测性能之间寻找帕累托最优。波束赋形中，通信波束需跟踪用户，感知波束需扫描目标区域，联合设计需要解决空间域资源的动态分配问题。

\textbf{（3）非合作目标的检测稳定性。}无人机的RCS随姿态角剧烈变化（可达20\,dB以上起伏），且小型多旋翼的RCS通常低于$-15$\,dBsm，导致在杂波环境中检测概率不稳定，虚警率和漏检率难以同时控制。

\subsection{未来演进方向}

ISAC技术的未来演进呈现三大趋势：

\textbf{（1）AI辅助感知。}利用深度学习替代传统CFAR检测器，如基于卷积神经网络（CNN）的Range-Doppler图目标检测，可在低信噪比条件下提升5--8\,dB的检测增益。同时，基于循环神经网络（RNN/LSTM）的轨迹预测可实现目标行为意图识别，从"检测-定位"的低层感知上升到"识别-理解"的语义感知层面。

\textbf{（2）RIS增强ISAC。}智能超表面（Reconfigurable Intelligent Surface, RIS）通过大量低成本反射单元实现对电磁波传播环境的可编程控制。在ISAC场景中，RIS可以：(a) 创建虚拟视距（Virtual LoS）路径，扩展非视距条件下的感知覆盖；(b) 增强目标回波信号强度，提升弱目标检测能力；(c) 形成空间分集，改善角度估计精度。

\textbf{（3）6G原生感知。}6G系统的目标是在空口设计之初就将感知作为原生能力纳入（Sensing-by-Design），而非5G-A阶段的"通信优先、感知叠加"模式。这意味着新的波形设计（如OTFS用于高多普勒场景）、新的帧结构（感知与通信时隙动态切换）和新的网络架构（感知即服务, Sensing-as-a-Service）。

\subsection{总结}

本报告以低空物联网为应用背景，对5G-A通感一体化技术进行了系统性的研讨分析。从OFDM雷达信号模型与2D-FFT算法的推导与仿真验证，到3GPP R18/R19/R20标准化进展的梳理，再到IoT平台端到端五层架构的设计和MQTT/CoAP/LwM2M协议选型分析，最后对三大运营商的低空试点和ISAC与传统雷达方案进行了对比评估，形成了"物理层感知$\to$数据上报$\to$平台处理$\to$应用驱动"的完整分析闭环。仿真结果表明，5G NR标准参数下的ISAC系统（$N_{\mathrm{CPI}}=14$时隙跨时隙处理）可实现1.5\,m距离分辨率、4.4\,m/s速度分辨率和5\,km探测距离，满足低空无人机监管需求。产业实践证明了ISAC方案的商业可行性。随着AI、RIS等使能技术的成熟和6G标准化的推进，通信感知一体化将成为未来移动通信网络的标配能力。
